\Needspace{12\baselineskip}
\begin{center}\textit{An analogous method for proving the theorem of Article 114.}\end{center}

\Article{145}
In the preceding proofs we always took an even value for $e$ (Articles 137\,..\,144). An odd value could also have been used, but we would then have had to introduce further distinctions. Those who enjoy these investigations will find it useful to try working out these cases for themselves. We would also then have had to assume the theorems concerning the residues $+2$ and $-2$. But since our proof was completed without these theorems, it gives us a new method of proving them. This is worth something, since the methods we used above to prove that $\pm2$ is a residue of every prime of the form $8n+1$ may seem less direct. We will assume that the remaining cases (concerning primes of the forms $8n+3$, $8n+5$, $8n+7$) have been proved by the methods given above, and that this theorem has only been found by examining examples. The following reasoning will turn this evidence into certainty.

If $\pm2$ were not a residue of all primes of the form $8n+1$, let the smallest prime of this form of which $\pm2$ is a nonresidue be $=a$, so that the theorem holds for all primes smaller than $a$. Then take a prime $<\frac{1}{2}a$ of which $a$ is a nonresidue (its existence follows easily from Article 129). Let it be $=p$. The fundamental theorem then gives $pNa$, and this gives $\pm2pRa$. --- So let $e^2\equiv2p\ (\mathrm{mod.}\ a)$, with $e$ odd and $<a$. We must now distinguish two cases.

I. When $e$ is not divisible by $p$, let $e^2=2p+aq$. Then $q$ is positive, of the form $8n+7$ or $8n+3$ (according to whether $p$ has the form $4n+1$ or $4n+3$), $<a$, and not divisible by $p$. Divide all the prime factors of $q$ into four classes: let there be $e$ of the form $8n+1$, $f$ of the form $8n+3$, $g$ of the form $8n+5$, and $h$ of the form $8n+7$. Let the product of the factors in the first class be $E$, and the products of the factors in the second, third and fourth classes respectively $F,G.H$\HistoricalNote{*)}{If a class contains no factors, we must write $1$ in place of their product.}.

\SourceLeaf{116}With this done, consider \textit{first} the case where $p$ has the form $4n+1$, or $q$ the form $8n+7$. It is easy to see that $2RE$, $2RH$, which gives $pRE$, $pRH$, and finally $ERp$, $HRp$. Also $2$ is a nonresidue of every factor of the form $8n+3$ or $8n+5$, and so is $p$; therefore each such factor is a nonresidue of $p$. It follows easily that $FG$ is a residue of $p$ if $f+g$ is even, and a nonresidue if $f+g$ is odd. But $f+g$ cannot be odd: listing all the cases shows that $EFGH$, or $q$, would have either the form $8n+3$ or $8n+5$ if $f+g$ were odd, whatever the individual $e,f,g,h$ might be, contrary to the hypothesis. Thus $FGRp$, $EFGHRp$, or $qRp$, and finally, because $aqRp$, $aRp$, contrary to the hypothesis. \textit{Second,} when $p$ has the form $4n+3$, we can similarly show that $pRE$, and therefore $ERp$; $-pRF$, and therefore $FRp$; and finally that $g+h$ is even, so $GHRp$. This finally gives $qRp$, $aRp$, contrary to the hypothesis.

II. When $e$ is divisible by $p$, the proof can be carried out in a similar way and worked out without difficulty by experienced readers (the only readers for whom this article is written). We omit it for brevity.

\Needspace{10\baselineskip}
\begin{center}\textit{Solution of the general problem.}\end{center}

\Article{146}
The fundamental theorem and the propositions concerning the residues $-1$ and $\pm2$ always let us determine whether any given number is a residue or nonresidue of a given prime. It will be useful, however, to bring together here the other results given above, so that we have everything needed to solve the following

\textsc{Problem:} \textit{Given any two numbers $P$, $Q$, determine whether the second, $Q$, is a residue or nonresidue of the first, $P$.}

\textit{Solution.} I. Let $P=a^{\alpha}b^{\beta}c^{\gamma}$, etc., where $a,b,c$, etc., denote distinct positive primes (since $P$ must clearly be taken in absolute value). For brevity, in this article we will call the \textit{relation} of two numbers $x,y$ simply the relation in which the first, $x$, is a residue or nonresidue of the second, $y$. Thus the relation of $Q,P$ depends on the relations of $Q,a^{\alpha}$; $Q,b^{\beta}$, etc. (Article 105).

II. To find the relation of $Q,a^{\alpha}$ (the same applies to the others, $Q,b^{\beta}$, etc.), we must distinguish two cases.

\SourceLeaf{117}1. When $Q$ is divisible by $a$, put $Q=Q'a^e$, so that $Q'$ is not divisible by $a$. If $e=\alpha$ or $e>\alpha$, then $QRa^{\alpha}$; if $e<\alpha$ and the exponent is odd, then $QNa^{\alpha}$; finally, if $e<\alpha$ and the exponent is even, $Q$ has the same relation to $a^{\alpha}$ as $Q'$ has to $a^{\alpha-e}$. This reduces the present case to

2. When $Q$ is not divisible by $a$. Here we again distinguish two cases.

$(A)$ When $a=2$. Then $QRa^{\alpha}$ always holds when $\alpha=1$. When $\alpha=2$, $Q$ must have the form $4n+1$; finally, when $\alpha=3$ or $>3$, $Q$ must have the form $8n+1$. If this condition holds, then $QRa^{\alpha}$.

$(B)$ When $a$ is another prime. Then $Q$ has the same relation to $a^{\alpha}$ as to $a$ (see Article 101).

III. The relation of any number $Q$ to an odd prime $a$ is found as follows. When $Q>a$, replace $Q$ by its least positive residue modulo $a$\HistoricalNote{*)}{\textit{Residue} in the sense of Article 4. --- It is usually better to take the \textit{absolutely} least residue.}. This has the same relation to $a$ as $Q$ does.

Next, resolve $Q$, or the number taken in its place, into its prime factors $p,p',p''$, etc., adding the factor $-1$ when $Q$ is negative. The relation of $Q$ to $a$ is known to depend on the relations of the individual $p,p',p''$, etc., to $a$. Specifically, if $2m$ of these factors are nonresidues of $a$, then $QRa$; if there are $2m+1$, then $QNa$. It is easy to see that if among the factors $p,p',p''$, etc., two, four, six, or generally $2k$ equal ones occur, they can safely be left out.

IV. If the factors $p.p',p''$ include $-1$ and $2$, their relation to $a$ can be found from Articles 108, 112, 113, 114. The relation of the remaining factors to $a$ depends on the relation of $a$ to them (the fundamental theorem and the propositions of Article 131). Let $p$ be one of them. Treating the numbers $a,p$ as we previously treated $Q$ and $a$, which are respectively larger, we find that the relation of $a$ to $p$ can either be determined from Articles 108---114 (if the least residue of $a$ $(\mathrm{mod.}\ p)$ has no odd prime factors), or depends in turn on the relation of $p$ to certain primes smaller than $p$. The same applies to the remaining factors $p',p''$, etc. It is now easy to see that \SourceLeaf{118}continuing this operation eventually brings us to numbers whose relations can be determined from the propositions of Articles 108---114. An example will make this clearer.

\textit{Example.} We seek the relation of $+453$ to $1236$. We have $1236=4.3.103$; $+453R4$ by II. 2 $(A)$; $+453R3$ by II. 1. So it remains to find the relation of $+453$ to $103$. This is the same as that of $+41$ $(\equiv453,\ \mathrm{mod.}\ 103)$ to $103$, and the same as that of $+103$ to $41$ (by the fundamental theorem), or of $-20$ to $41$. But $-20R41$, since $-20=-1.2.2.5$; $-1R41$ (Article 108); and $+5R41$, because $41\equiv1$ and it is therefore a residue of $5$ (by the fundamental theorem). Thus $+453R103$, and finally $+453R1236$. In fact, $453\equiv297^2\ (\mathrm{mod.}\ 1236)$.

\Needspace{12\baselineskip}
\begin{center}\textit{Linear forms containing all primes of which any given number is a residue or nonresidue.}\end{center}

\Article{147}
Given any number $A$, we can exhibit certain \textit{formulas} containing all numbers coprime to $A$ of which $A$ is a residue, or all numbers that can be \textit{divisors} of numbers of the form $xx-A$ (where $xx$ denotes an indeterminate square)\HistoricalNote{*)}{We will simply call such numbers \textit{divisors of $xx-A$}; the meaning of \textit{non-divisors} is then clear.}. For brevity, however, we consider only those divisors that are odd and coprime to $A$, since the remaining cases easily reduce to these.

First, let $A$ be either a positive prime of the form $4n+1$ or a negative prime of the form $4n-1$. By the fundamental theorem, all primes that, taken positively, are residues of $A$ will be divisors of $xx-A$; all primes (except the number $2$, which is always a divisor) that are nonresidues of $A$ will be non-divisors of $xx-A$. Let all residues of $A$ smaller than $A$ (excluding zero) be $r,r',r''$, etc., and all nonresidues $n,n',n''$, etc. Every prime contained in one of the forms $Ak+r$, $Ak+r'$, $Ak+r''$, etc., will be a divisor of $xx-A$, while every prime in one of the forms $Ak+n$, $Ak+n'$, etc., will be a non-divisor, where $k$ denotes an indeterminate integer. We call the former \textit{divisor forms of $xx-A$}, and the latter \textit{non-divisor forms}. There are $\frac{1}{2}(A-1)$ of each. Further, if $B$ is an odd composite number and $ARB$, all prime factors of $B$ lie in one of the \SourceLeaf{119}former forms, and so does $B$. Thus \textit{every} odd number contained in a non-divisor form is a non-divisor of $xx-A$. But this theorem cannot be reversed: if $B$ is an odd composite non-divisor of $xx-A$, some prime factors of $B$ will be non-divisors; if their number is \textit{even}, $B$ will nevertheless lie in a divisor form. See Article 99.

\textit{Example.} In this way, for $A=-11$ we find the following divisor forms of $xx+11$: $11k+\nobreak{}1,\allowbreak{}3,\allowbreak{}4,\allowbreak{}5,\allowbreak{}9$; the non-divisor forms are $11k+\nobreak{}2,\allowbreak{}6,\allowbreak{}7,\allowbreak{}8,\allowbreak{}10$. Thus $-11$ is a nonresidue of all odd numbers contained in one of the latter forms, but a residue of all primes belonging to one of the former.

\Needspace{12\baselineskip}
Similar forms exist for the divisors and non-divisors of $xx-A$, whatever number $A$ denotes. It is easy to see, however, that we need consider only values of $A$ not divisible by any square. For if $A=a^2A'$, all divisors\HistoricalNote{*)}{Those coprime to $A$, of course.} of $xx-A$ are also divisors of $xx-A'$, and the same holds for non-divisors. --- We distinguish three cases: 1) when $A$ has the form $+(4n+1)$ or $-(4n-1)$; 2) when $A$ has the form $-(4n+1)$ or $+(4n-1)$; 3) when $A$ is even, or has the form $\pm(4n+2)$.

\Needspace{10\baselineskip}
\Article{148}
\textit{First case:} when $A$ has the form $+(4n+1)$ or $-(4n-1)$. Resolve $A$ into its prime factors, giving a positive sign to those of the form $4n+1$, and a negative sign to those of the form $4n-1$ (so their product is $=A$). Let these factors be $a,b,c,d$, etc. Divide all numbers smaller than $A$ and coprime to $A$ into two classes. Put in the first class all numbers that are nonresidues of none of $a,b,c,d$, etc., or of two, or four, or generally an even number of them; in the second, those that are nonresidues of one of $a,b,c$, etc., or of three, etc., or generally an odd number. Denote the former by $r,r',r''$, etc., and the latter by $n,n',n''$, etc. Then $Ak+r$, $Ak+r'$, $Ak+r''$, etc., are divisor forms of $xx-A$, while $Ak+n$, $Ak+n'$, etc., are non-divisor forms of $xx-A$ (\textit{that is, every prime except $2$ is a divisor or non-divisor of $xx-A$ according to whether it lies in one of the former or latter forms}). For if $p$ is \SourceLeaf{120}a positive prime and a residue or nonresidue of one of $a,b,c$, etc., that number itself is a residue or nonresidue of $p$ (by the fundamental theorem). So if $a,b,c$, etc., include $m$ numbers of which $p$ is a nonresidue, just as many are nonresidues of $p$. Thus if $p$ lies in one of the former forms, $m$ is even and $ARp$; if it lies in one of the latter, $m$ is odd and $ANp$.

\textit{Example.} Let $A=+105=-3\times+5\times-7$. Then $r,r',r''$, etc., are the following: $1,\allowbreak{}4,\allowbreak{}16,\allowbreak{}46,\allowbreak{}64,\allowbreak{}79$ (nonresidues of none of $3,\allowbreak{}5,\allowbreak{}7$); $2,\allowbreak{}8,\allowbreak{}23,\allowbreak{}32,\allowbreak{}53,\allowbreak{}92$ (nonresidues of $3,\allowbreak{}5$); $26,\allowbreak{}41,\allowbreak{}59,\allowbreak{}89,\allowbreak{}101,\allowbreak{}104$ (nonresidues of $3,\allowbreak{}7$); $13,\allowbreak{}52,\allowbreak{}73,\allowbreak{}82,\allowbreak{}97,\allowbreak{}103$ (nonresidues of $5,\allowbreak{}7$). --- The numbers $n,n',n''$, etc., are the following: $11,\allowbreak{}29,\allowbreak{}44,\allowbreak{}71,\allowbreak{}74,\allowbreak{}86$; $22,\allowbreak{}37,\allowbreak{}43,\allowbreak{}58,\allowbreak{}67,\allowbreak{}88$; $19,\allowbreak{}31,\allowbreak{}34,\allowbreak{}61,\allowbreak{}76,\allowbreak{}94$; $17,\allowbreak{}38.47,\allowbreak{}62,\allowbreak{}68,\allowbreak{}83$. The first six are nonresidues of $3$, the next six nonresidues of $5$, followed by the nonresidues of $7$, and finally those that are nonresidues of all three at once.

The theory of combinations and Articles 32, 96 easily show that the number of $r,r',r''$, etc., is
\[
=t\left(1+\frac{l\, .\, l-1}{1\, .\, 2}+\frac{l\, .\, l-1\, .\, l-2\, .\, l-3}{1\, .\, 2\, .\, 3\, .\, 4}+ .\, .\, .\, .\, .\right)
\]
and the number of $n,n',n''$, etc.,
\[
=t\left(l+\frac{l\, .\, l-1\, .\, l-2}{1\, .\, 2\, .\, 3}+\frac{l\, .\, l-1\, .\, .\, .\, .\, l-4}{1\, .\, 2\, .\, .\, .\, .\, 5}+ .\, .\, .\right)
\]
where $l$ denotes the number of $a,b,c$, etc.;
\[
t=2^{-l}(a-1)(b-1)(c-1)\text{ etc.}
\]
and each series is continued until it terminates. (There are $t$ numbers that are residues of all of $a,b,c$, etc., $\frac{t\, .\, l\, .\, l-1}{1\, .\, 2}$ that are nonresidues of two, etc.; but brevity does not permit us to explain this proof more fully.) The sum of each series\HistoricalNote{*)}{Leaving out the factor $t$.} is $=2^{l-1}$. The first comes from
\[
1+(l-1)+\frac{l-1\, .\, l-2}{1\, .\, 2}+ .\, .\, .
\]
by joining the second and third terms, the fourth and fifth, etc.; the second comes from the same series by joining the first and second terms, the third and fourth, etc. Thus there are as many divisor forms of $xx-A$ as non-divisor forms, namely $\frac{1}{2}(a-1)(b-1)(c-1)$, etc.

\SourceLeaf{121}\Needspace{10\baselineskip}
\Article{149}
We can consider the \textit{second and third cases} together here. We can always put $A$ $=(-1)Q$, or $=(+2)Q$, or $=(-2)Q$, where $Q$ denotes a number of the form $+(4n+1)$ or $-(4n-1)$, such as we considered in the preceding article. In general, let $A=\alpha Q$, with $\alpha$ either $=-1$ or $=\pm2$. Then $A$ is a residue of all numbers of which either both $\alpha$ and $Q$ are residues or neither is; it is a nonresidue of all numbers of which exactly one of $\alpha,Q$ is a nonresidue. The divisor and non-divisor forms of $xx-A$ follow easily. If $\alpha=-1$, divide all numbers smaller than $4A$ and coprime to it into two classes. Put in the first those lying in a divisor form of $xx-Q$ and also in the form $4n+1$, and those lying in a non-divisor form of $xx-Q$ and also in the form $4n+3$; put the rest in the second. Let the former be $r,r',r''$, etc., and the latter $n,n',n''$, etc. Then $A$ is a residue of all primes in one of the forms $4Ak+r$, $4Ak+r'$, $4Ak+r''$, etc., and a nonresidue of all primes in one of the forms $4Ak+n$, $4Ak+n'$, etc. --- If $\alpha=\pm2$, divide all numbers smaller than $8Q$ and coprime to it into two classes. Put in the first those lying in a divisor form of $xx-Q$ and also in one of the forms $8n+1$, $8n+7$ for the upper sign, or $8n+1$, $8n+3$ for the lower sign; also those lying in a non-divisor form of $xx-Q$ and in one of $8n+3$, $8n+5$ for the upper sign, or $8n+5$, $8n+7$ for the lower sign. Put the rest in the second class. Denoting the numbers of the first class by $r,r',r''$, etc., and those of the second by $n,n',n''$, etc., $\pm2Q$ is a residue of all primes in one of the forms $8Qk+r$, $8Qk+r'$, $8Qk+r''$, etc., and a nonresidue of all primes in one of the forms $8Qk+n$, $8Qk+n'$, $8Qk+n''$, etc. It can easily be proved that here too there are as many divisor forms of $xx-A$ as non-divisor forms.

\textit{Example.} This gives $+10$ as a residue of all primes in one of the forms $40k+\nobreak{}1,\allowbreak{}3,\allowbreak{}9,\allowbreak{}13,\allowbreak{}27,\allowbreak{}31,\allowbreak{}37,\allowbreak{}39$, but a nonresidue of all primes contained in one of $40k+\nobreak{}7,\allowbreak{}11,\allowbreak{}17,\allowbreak{}19,\allowbreak{}21,\allowbreak{}23,\allowbreak{}29,\allowbreak{}33$.

\SourceLeaf{122}\Needspace{8\baselineskip}
\Article{150}
These forms have several quite notable properties, but we give just one. If $B$ is a composite number coprime to $A$, and $2m$ of its prime factors lie in a non-divisor form of $xx-A$, then $B$ lies in a divisor form of $xx-A$. If, however, the number of prime factors of $B$ lying in a non-divisor form of $xx-A$ is odd, $B$ too lies in a non-divisor form. We omit the proof, which is not difficult. It follows that every prime, and also every odd composite number coprime to $A$ that lies in a non-divisor form, is a non-divisor: some prime factor of such a number must necessarily be a non-divisor.

\Needspace{12\baselineskip}
\begin{center}\textit{The work of others on these investigations.}\end{center}

\Article{151}
The fundamental theorem, certainly one of the most elegant results of this kind, has not so far been stated by anyone in the simple form given above. This is all the more surprising since the distinguished Euler already knew certain other propositions resting on it, from which he could easily have returned to it. He knew that there are certain forms containing all prime divisors of numbers of the form $xx-A$, and other forms containing all prime non-divisors of numbers of that same form, so that the latter exclude the former; and he had found a method for determining these forms. But all his attempts to reach a proof were unsuccessful and only made this truth, found from examples, more probable. In his paper \textit{Novae demonstrationes circa divisores numerorum formae $xx+nyy$}, read to the St Petersburg Academy on 20 November 1775 and preserved after the death of that great man in \textit{Vol.} I of the Academy's \textit{Nov. Act.}, p. 47 ff., he does seem to have thought he had succeeded. But an error entered here: on p. 65 he \textit{tacitly} assumed that such divisor and non-divisor forms exist\HistoricalNote{*}{Namely, that there are numbers $r,r',r''$, etc.; $n,n',n''$, etc., all different and $<4A$, such that all prime divisors of $xx-A$ lie in one of the forms $4Ak+r$, $4Ak+r'$, etc., and all prime non-divisors in one of $4Ak+n$, $4Ak+n'$, etc. (where $k$ denotes an indeterminate number).}. From this it was not difficult to derive \textit{what} they must be, but the method he used to support this supposition does not seem \SourceLeaf{123}suitable. In another paper, \textit{De criteriis aequationis $fxx+gyy=hzz$ utrumque resolutionem admittat necne}, Opusc. Anal., Vol. I (where $f,g,h$ are given, and $x,y,z$ indeterminate), he found by examining examples that if the equation is solvable for some value $h=s$, it is also solvable for every other value congruent to $s$ modulo $4fg$, provided it is a prime. The supposition we mentioned can easily be proved from this proposition. But the proof of this theorem too eluded all his efforts\HistoricalNote{*)}{As he admits himself, in the cited work, p. 216: ``A proof of this most elegant theorem is still wanted, after many have long sought it in vain.\,\ldots\ Whoever succeeds in finding a proof of this theorem must therefore be thought to have achieved a great deal.'' --- Many other passages of Opusc. Anall. show how eagerly that immortal man wanted a proof of this theorem and others that are only special cases of the fundamental theorem. Compare \textit{Additamentum ad diss.} VIII, Vol. I, and \textit{diss.} XIII, Vol. II, and several papers in Comment. Petrop., already cited in various places.}. This is not surprising, since in our view one had to start from the fundamental theorem. The truth of this proposition will, moreover, follow directly from what we teach in the next section.

After Euler, the distinguished Le Gendre worked on the same subject in his excellent paper \textit{Recherches d'analyse indéterminée}, Hist. de l'Ac. des Sc., 1785, p. 465 ff. He reached a theorem equivalent in substance to the fundamental theorem: if $p,q$ denote two positive primes, the absolutely least residues of the powers $p^{\frac{q-1}{2}}$, $q^{\frac{p-1}{2}}$ modulo $q,p$, respectively, are either both $+1$ or both $-1$ when either $p$ or $q$ has the form $4n+1$. When both $p$ and $q$ have the form $4n+3$, one least residue is $+1$ and the other $-1$ (p. 516). By Article 106, it follows that the \textit{relation} (in the sense of Article 146) of $p$ to $q$ and of $q$ to $p$ is \textit{the same} when either $p$ or $q$ has the form $4n+1$, and \textit{opposite} when both $p$ and $q$ have the form $4n+3$. This proposition is among those of Article 131; it also follows from 1, 3, 9 of Article 133. Conversely, the fundamental theorem can be derived from it. Le Gendre also attempted a proof; since it is highly ingenious, we will discuss it more fully in the next section. But he assumed several things without proof (as he himself admits on p. 520, \textit{We have only assumed}, etc.), some of which have so far been proved by no one, and some, in our view at least, cannot be proved without the fundamental theorem itself. So his approach seems unable to reach the goal, and our proof must be regarded as the first. --- We will, moreover, give \textit{two other proofs} of this most important theorem below, entirely different from the preceding proof and from each other.

\SourceLeaf{124}\Needspace{12\baselineskip}
\begin{center}\textit{Quadratic congruences with a linear term.}\end{center}

\Article{152}
So far we have treated the pure congruence $xx\equiv A\ (\mathrm{mod.}\,m)$ and shown how to decide whether it is solvable. Finding \textit{the roots themselves} was reduced in Article 105 to the case where $m$ is a prime or a power of a prime, and the latter case in Article 101 to the case where $m$ is a prime. For this case, what we gave in Article 61 ff., together with what we will teach in Sections V and VIII, contains almost everything that can be obtained by direct methods. But where these methods apply, they are usually infinitely longer than the indirect methods we will teach in Section VI, and are therefore notable more for their beauty than for their practical usefulness. --- \textit{Quadratic congruences with a linear term} can easily be reduced to pure ones. Given the congruence
\[
axx+bx+c\equiv 0
\]
to be solved modulo $m$, the following congruence is equivalent to it:
\[
4aaxx+4abx+4ac\equiv 0\ (\mathrm{mod.}\,4am)
\]
\textit{that is,} every number satisfying either one also satisfies the other. This can be written as
\[
(2ax+b)^2\equiv bb-4ac\ (\mathrm{mod.}\,4am)
\]
from which all values of $2ax+b$ smaller than $4am$, if there are any, can be found. Denoting these by $r,r',r''$, etc., all solutions of the proposed congruence follow from the solutions of the congruences
\[
2ax\equiv r-b,\quad 2ax\equiv r'-b\text{ etc.}\ (\mathrm{mod.}\,4am)
\]
which we taught how to find in Section II. We also observe that various tricks can usually shorten the solution. \textit{For example,} in place of the proposed congruence one can find an equivalent congruence
\[
a'xx+2b'x+c'\equiv 0
\]
in which $a'$ divides $m$. Brevity does not allow us to explain this here; the last section may be consulted on these matters.
